\documentclass[11pt]{article}

\usepackage[margin=0.75in]{geometry}
\usepackage{amsmath,amssymb,booktabs,xcolor,enumitem}
\usepackage[colorlinks=true,linkcolor=blue,citecolor=blue,urlcolor=blue]{hyperref}

\newcommand{\sg}{\operatorname{sg}}
\newcommand{\E}{\mathbb{E}}
\newcommand{\inner}[2]{\langle #1,\, #2\rangle}
\newcommand{\Fmap}{F}
\newcommand{\Vmap}{F_V}
\newcommand{\Bmap}{B}
\newcommand{\Lexp}{L^{\tau}_2}
\newcommand{\ind}{\mathbf{1}}
\newcommand{\off}{\operatorname{off}}
\newcommand{\clip}{\operatorname{clip}}

\title{\textbf{In-Sample Forward--Backward}\\[3pt]\large Adding IQL-style improvement to successor-measure learning}
\author{}
\date{}

\begin{document}
\maketitle

\section{Core idea}

Forward--Backward learns a successor measure.

\begin{equation}
M(s,a,s';z)=\frac{1}{\sqrt d}\inner{\Fmap(s,a,z)}{\Bmap(s')}.
\end{equation}

This quantity says how much future occupancy of state $s'$ is predicted after taking action $a$ in state $s$ under task direction $z$.
The actor uses the scalar projection of this measure onto the task direction.

\begin{equation}
Q_z(s,a)=\frac{1}{\sqrt d}\inner{\Fmap(s,a,z)}{z}.
\end{equation}

For goal reaching, one typically sets $z=\Bmap(g)$.
The actor then chooses actions whose predicted future occupancy points toward the goal embedding.
Standard FB trains this measure with a policy-evaluation backup.

\begin{equation}
M(s,a,\cdot;z)\leftarrow \delta_{s'}+\gamma M(s',a^+,\cdot;z),\qquad a^+\sim\pi_z(\cdot\mid s').
\end{equation}

This asks FB to evaluate the future occupancy of the current policy.
It does not explicitly ask which dataset-supported actions are better for the task.
The bridge experiments suggest that this missing improvement operator is the central difference between FB and IQL.

\begin{center}
\begin{tabular}{@{}lll r@{}}
\toprule
\textbf{Configuration} & \textbf{Objective} & \textbf{Architecture} & \textbf{Success} \\
\midrule
GCIQL bilinear base & IQL & bilinear $\phi^\top\psi/\sqrt d$ & ${\sim}90\%$ \\
B4 & IQL & full FB skeleton & ${\sim}87\%$ \\
C4 & measure & full FB skeleton & ${\sim}16\%$ \\
Measure variants & measure & bilinear/FB variants & ${\sim}9$--$25\%$ \\
\bottomrule
\end{tabular}
\end{center}

The bilinear architecture works well under the IQL objective.
The point is that IQL performs in-sample policy improvement, while standard FB mostly performs successor-measure evaluation.

\section{The addition}

The addition is a state-only value successor measure.

\begin{equation}
N(s,s';z)=\frac{1}{\sqrt d}\inner{\Vmap(s,z)}{\Bmap(s')}.
\end{equation}

This mirrors the relationship between $Q$ and $V$ in IQL.

\begin{center}
\begin{tabular}{@{}lll@{}}
\toprule
\textbf{Role} & \textbf{IQL} & \textbf{ISFB} \\
\midrule
State-action object & $Q(s,a,g)$ & $M(s,a,\cdot;z)$ \\
State value object & $V(s,g)$ & $N(s,\cdot;z)$ \\
Improvement step & $V$ favors better dataset actions & $N$ favors better dataset actions \\
Bootstrap & $Q\leftarrow r+\gamma V'$ & $M\leftarrow \delta_{s'}+\gamma N'$ \\
\bottomrule
\end{tabular}
\end{center}

The new FB backup replaces the sampled next action with the value successor measure.

\begin{equation}
M(s,a,\cdot;z)\leftarrow \delta_{s'}+\gamma N(s',\cdot;z).
\end{equation}

So standard FB asks what future follows the current actor.
ISFB asks what future follows after an improved in-sample choice.
For a minibatch $\{(s_i,a_i,s'_i)\}_{i=1}^B$, define the action and value measures as follows.

\begin{equation}
M_{ij}=\frac{1}{\sqrt d}\inner{\Fmap(s_i,a_i,z_i)}{\Bmap(s'_j)},\qquad N_{ij}=\frac{1}{\sqrt d}\inner{\Vmap(s_i,z_i)}{\Bmap(s'_j)}.
\end{equation}

The three versions differ only in how they make $N$ represent an improved state-level successor measure.

\section{Version 1: coordinatewise ISFB}

Coordinatewise ISFB applies the expectile directly to each successor coordinate.

\begin{equation}
\mathcal L_N^{\mathrm{coord}}=\frac{1}{B^2}\sum_{i,j}\Lexp\!\left(\sg M_{ij}-N_{ij}\right).
\end{equation}

This is the most direct measure-level analogue of IQL.

\begin{equation}
V(s)\approx \operatorname{Expectile}_\tau Q(s,a)\quad\Longrightarrow\quad N(s,s';z)\approx \operatorname{Expectile}_\tau M(s,a,s';z).
\end{equation}

The issue is that the improvement decision is made separately for every successor coordinate.
Different coordinates may prefer different actions.
Thus $N(s,\cdot;z)$ can become an optimistic vector that no single action or policy can realize.
This makes coordinatewise ISFB a useful baseline, but not the cleanest policy-improvement object.

\section{Version 2: scalar-gated ISFB}

Scalar-gated ISFB first makes the improvement decision in the scalar value space used by the actor.

\begin{equation}
q_i=Q_z(s_i,a_i)=\frac{1}{\sqrt d}\inner{\Fmap(s_i,a_i,z_i)}{z_i}.
\end{equation}

\begin{equation}
v_i=V_z(s_i)=\frac{1}{\sqrt d}\inner{\Vmap(s_i,z_i)}{z_i}.
\end{equation}

The scalar value leg is trained with the usual IQL expectile.

\begin{equation}
\mathcal L_V^{\mathrm{scalar}}=\frac{1}{B}\sum_i\Lexp\!\left(\sg q_i-v_i\right).
\end{equation}

The scalar advantage is then used as the action-selection signal.

\begin{equation}
A_i=\sg(q_i-v_i).
\end{equation}

The hard scalar gate uses the expectile weights directly.

\begin{equation}
w_i^{\mathrm{hard}}=\tau\ind[A_i>0]+(1-\tau)\ind[A_i\le 0].
\end{equation}

The whole successor vector is then copied strongly or weakly depending on whether the action looks good for task $z_i$.

\begin{equation}
\mathcal L_N^{\mathrm{hard}}=\frac{1}{B^2}\sum_{i,j}\bar w_i^{\mathrm{hard}}\left(\sg M_{ij}-N_{ij}\right)^2,\qquad \bar w_i^{\mathrm{hard}}=\frac{w_i^{\mathrm{hard}}}{\frac{1}{B}\sum_k w_k^{\mathrm{hard}}+\epsilon}.
\end{equation}

The key difference from coordinatewise ISFB is that action selection happens once per transition, after projecting onto the task direction.
If an action is good for the task, its entire successor measure is used.
This is more policy-coherent than choosing winners independently for each future coordinate.

\section{Version 3: AWR-gated ISFB}

AWR-gated ISFB uses the same scalar advantage but replaces the hard gate with a smooth exponential weight.

\begin{equation}
w_i^{\mathrm{AWR}}=\clip\!\left(\exp(\beta A_i),0,w_{\max}\right).
\end{equation}

The value-measure regression is then weighted by this smooth advantage weight.

\begin{equation}
\mathcal L_N^{\mathrm{AWR}}=\frac{1}{B^2}\sum_{i,j}\bar w_i^{\mathrm{AWR}}\left(\sg M_{ij}-N_{ij}\right)^2,\qquad \bar w_i^{\mathrm{AWR}}=\frac{w_i^{\mathrm{AWR}}}{\frac{1}{B}\sum_k w_k^{\mathrm{AWR}}+\epsilon}.
\end{equation}

The hard gate only uses the sign of the advantage.
AWR also uses the magnitude of the advantage.
This can make the improvement step sharper.
It also makes the method sensitive to the scale of $A_i$.
If advantages are too small, all weights become nearly uniform.
If advantages are too large, a few samples dominate after clipping.

\section{Shared ISFB measure TD}

All three versions use the same measure TD update for $M$.

\begin{equation}
N^{-}_{ij}=\frac{1}{\sqrt d}\inner{\Vmap^{-}(s'_i,z_i)}{\Bmap^{-}(s'_j)}.
\end{equation}

\begin{equation}
\Delta_{ij}^{\mathrm{ISFB}}=M_{ij}-\gamma\sg N^{-}_{ij}.
\end{equation}

With a termination or goal-success mask $m_i$, the residual becomes masked row-wise.

\begin{equation}
\Delta_{ij}^{\mathrm{ISFB}}=M_{ij}-\gamma m_i\sg N^{-}_{ij}.
\end{equation}

The measure loss keeps the FB off-diagonal and diagonal structure.

\begin{equation}
\mathcal L_M^{\mathrm{ISFB}}=\frac{1}{2|\off|}\sum_{i\ne j}\left(\Delta_{ij}^{\mathrm{ISFB}}\right)^2-\frac{1}{B}\sum_i\Delta_{ii}^{\mathrm{ISFB}}.
\end{equation}

The backward embedding keeps the standard FB orthonormality loss.

\begin{equation}
\mathcal L_\perp=\frac{1}{2|\off|}\sum_{i\ne j}\inner{\Bmap(s'_i)}{\Bmap(s'_j)}^2-\frac{1}{B}\sum_i\|\Bmap(s'_i)\|^2.
\end{equation}

\section{Complete objectives}

The coordinatewise objective is as follows.

\begin{equation}
\boxed{\mathcal L_{\mathrm{Coord\text{-}ISFB}}=\alpha_N\mathcal L_N^{\mathrm{coord}}+\alpha_M\mathcal L_M^{\mathrm{ISFB}}+\lambda\mathcal L_\perp.}
\end{equation}

The hard scalar-gated objective is as follows.

\begin{equation}
\boxed{\mathcal L_{\mathrm{Hard\text{-}ISFB}}=\mathcal L_V^{\mathrm{scalar}}+\alpha_N\mathcal L_N^{\mathrm{hard}}+\alpha_M\mathcal L_M^{\mathrm{ISFB}}+\lambda\mathcal L_\perp.}
\end{equation}

The AWR-gated objective is as follows.

\begin{equation}
\boxed{\mathcal L_{\mathrm{AWR\text{-}ISFB}}=\mathcal L_V^{\mathrm{scalar}}+\alpha_N\mathcal L_N^{\mathrm{AWR}}+\alpha_M\mathcal L_M^{\mathrm{ISFB}}+\lambda\mathcal L_\perp.}
\end{equation}

We should begin with the FB-scale orthonormality coefficient, such as $\lambda=1$, rather than large cube-scale values such as $\lambda=1000$.

% ============================================================================
\section*{Measured result: FB vs.\ SG-ISFB on cube ($\lambda=1$)}
% ============================================================================

Cube-single state, flow (FQL) actor throughout, orthonormality $\lambda=1$ for
every row, 3 seeds @ 500k (torch FB agent; wandb \texttt{amsks/factored-fb}).

\begin{center}
\begin{tabular}{@{}l l r@{}}
\toprule
\textbf{Config} (ortho $\lambda{=}1$, flow actor) & \textbf{Objective} & \textbf{Succ.} \\
\midrule
FB $+$ flow (C4: z-fwd $+$ sphere)   & measure, no in-sample operator          & \bad{16.5} \\
\midrule
ISFB $+$ flow, coord                 & measure $+$ coordinatewise expectile    & \good{33.3} \\
ISFB $+$ flow, scalar-gate AWR       & measure $+$ scalar gate, $e^{\beta A}$   & 26.7 \\
ISFB $+$ flow, scalar-gate hard      & measure $+$ scalar gate, hard expectile & \bad{8.0} \\
\bottomrule
\end{tabular}
\end{center}

\noindent\footnotesize
The FB baseline is the $\lambda{=}1$ FB-measure config (C4); the plain-FB anchor
itself uses $\lambda{=}1000$ ($\sim$48\% @1M). At matched ortho, the in-sample
operator roughly doubles the FB-measure floor ($16.5\to33.3$, coord) but stays far
below the IQL band ($\sim$87). The scalar-gate refinement underperforms the
coordinatewise baseline ($33.3\to26.7$ AWR $/\,8.0$ hard).
\normalsize

\end{document}